\begin{abstract}
본 보고서는 Glot500의 horizontal scaling 방향을 작은 재현 가능한 설정으로 가져와,
XLM-R 기반 multilingual encoder를 low-resource tail language로 확장할 때 새 vocabulary
embedding row의 초기화가 얼마나 중요한지 분석한다. 실험은 92개의 XLM-R-seen
language-script와 downstream 평가가 가능한 7개의 XLM-R-unseen target language-script를
사용한다. Tokenizer는 Glot500과 같이 mixed corpus에서 SentencePiece unigram tokenizer를
학습하고 기존 XLM-R SentencePiece model 뒤에 새 piece를 append하는 방식으로 고정한다.
따라서 main ablation은 tokenizer 자체가 아니라 새 token embedding row의 initialization,
즉 \method{random}, \method{mean}, \method{FVT}, \method{weighted FVT},
\method{family-aware mean}의 차이다.

최종 보고서는 50K-step convergence run을 중심으로 다섯 초기화 방법의 training loss,
pseudoperplexity, downstream scores를 tail/head/all group으로 정리한다. Step-4000 결과는
수렴 결과가 아니라 early diagnostic evidence로만 사용한다. 이 구분은 중요하다. 4K 이하
checkpoint에서는 \method{FVT}가 target pseudoperplexity와 여러 retrieval/alignment task에서
강한 초기 신호를 보였지만, 최종 주장은 50K run에서 loss flattening과 downstream trajectory가
확인된 뒤에만 세운다. 따라서 본 보고서는 특정 initialization이 모든 task에서 항상 최고라고
주장하지 않는다. 대신 같은 tokenizer, corpus, MLM objective, evaluation protocol을 공유하는
controlled setup에서 새 token embedding initialization이 low-resource vocabulary adaptation의
수렴 양상과 평가 지표에 실질적인 영향을 주는지를 검증한다.
\end{abstract}
